\section{Discussion}\label{sec:discussion}

\subsection{Exponential versus power-law convergence}

A central finding of this work is that convergence of CoT iterations on finite
state spaces is \emph{always exponential}, governed by the spectral gap of the
transition operator. The appearance of power-law convergence, as suggested by the
hitting time result $\widetilde{\Theta}(KM/\varepsilon)$ of Kim et al.\
\cite{kim2025metastable}, arises from two mechanisms.

First, in the metastable regime with small $\varepsilon$, the convergence rate
$1 - K\varepsilon/(K-1) \approx 1$ is close to unity. On any practical time
horizon, this slow exponential is indistinguishable from a power law. The Kim
et al.\ hitting time characterizes first passage to a \emph{specific} target
state rather than distributional convergence, and is consistent with our
spectral gap result since the mixing time scales as
$O((K-1)/(K\varepsilon) \cdot \log(1/\delta))$.

Second, the separation between within-cluster and between-cluster timescales
creates multi-scale dynamics: fast initial convergence (within-cluster
equilibration) followed by slow convergence (between-cluster mixing). This
two-phase behavior mimics power-law decay on intermediate timescales. Our
computational experiments confirm that exponential fits achieve $R^2 > 0.99$
across the entire range of $\varepsilon$, while power-law fits are consistently
inferior.

\subsection{The spectral gap as the fundamental quantity}

Theorems~\ref{thm:spectral}--\ref{thm:metastable} establish the spectral gap
$\gamma$ as the fundamental quantity governing convergence. The relationship
between $\gamma$ and convergence rate takes two forms:
\begin{itemize}[leftmargin=*]
\item In total variation: rate $= 1 - \gamma$ (asymptotically exact).
\item In $W_2$: rate $\approx (1 - \gamma)^{1/2}$
  (verified empirically; see Remark~\ref{rem:W2_rate}).
\end{itemize}
The square-root relationship in $W_2$ is a consequence of the comparison
inequality between $W_2$ and total variation on finite metric spaces
(Remark~\ref{rem:TV_W}). It remains open whether this bound is tight; see
Section~\ref{sec:open}.

\subsection{Connections to prior work}

The hitting time result of Kim et al.\ \cite{kim2025metastable},
$\widetilde{\Theta}(KM/\varepsilon)$, is consistent with our spectral gap
$\gamma = K\varepsilon/(K-1)$: the mixing time is $O(1/\gamma)$, and the
factor of $M$ in their result reflects that hitting a specific state within a
cluster of size $M$ incurs an additional $\Theta(M)$ factor beyond
distributional mixing.

The exponential convergence rate $K^L$ for looped networks established by
Ke et al.\ \cite{ke2024fixed} corresponds to
Theorem~\ref{thm:exponential}, where $K < 1$ plays the role of the Dobrushin
or Wasserstein contraction coefficient.

The JKO convergence results of Cheng et al.\ \cite{cheng2024convergence} for
flow-based generative models in $W_2$ under $\lambda$-convexity provide a
continuous-space parallel to our Theorem~\ref{thm:spectral}. Their exponential
convergence in $W_2$ via proximal gradient descent in Wasserstein space
corresponds to our discrete spectral convergence, with $\lambda$-convexity
playing the role of the spectral gap.

The eigenvalue classification of convergence regimes by Shukla and Joshi
\cite{shukla2025sde}---exponential, oscillatory, and boundary
attraction---parallels our spectral characterization adapted from continuous
SDEs to discrete Markov operators.

\subsection{Implications for CoT prompting}

Our results yield several practical implications.

\emph{Number of reasoning steps.} The required number of CoT steps scales as
$O(1/\gamma \cdot \log(1/\delta))$ for accuracy $\delta$. For metastable
problems with $K$ reasoning stages and inter-stage coupling $\varepsilon$,
this is $O((K-1)/(K\varepsilon) \cdot \log(1/\delta))$.

\emph{Diminishing returns.} Due to exponential convergence, the marginal
benefit of each additional CoT step decreases geometrically. This suggests the
existence of an optimal stopping time beyond which additional reasoning steps
provide negligible improvement.

\emph{Multiple solutions.} When the transition operator has multiple
near-fixed-points (metastable clusters), the CoT process may converge to an
incorrect fixed point depending on initialization. This provides a mathematical
explanation for why prompt engineering matters: the initial prompt directs the
reasoning trajectory toward the basin of attraction of the correct solution.

\subsection{Open questions}\label{sec:open}

We highlight several directions for future research.

\begin{conjecture}[Tight $W_2$ contraction]\label{conj:W2}
There exists a direct contraction bound
$W_2(T\mu, T\nu) \le \kappa \cdot W_2(\mu, \nu)$ for $\kappa < 1$ under
appropriate conditions on $P$, without passing through total variation. This
would yield a convergence rate of $\kappa$ rather than $\sqrt{\kappa}$ in
$W_2$.
\end{conjecture}

A direct $W_2$ contraction would sharpen Theorem~\ref{thm:spectral}(ii)
considerably. The obstacle is that the Kantorovich--Rubinstein duality, which
underlies our $W_1$ contraction proof, does not extend directly to $W_2$.

\begin{conjecture}[Non-Markovian power laws]\label{conj:nonmarkov}
For CoT processes with memory (non-Markovian dynamics), genuine power-law
convergence $W_2(T^n \mu_0, \mu^*) = \Theta(n^{-\alpha})$ can emerge,
with exponent $\alpha$ depending on the memory kernel.
\end{conjecture}

All Markov chains on finite state spaces converge exponentially. Genuine
power-law convergence would require infinite state spaces or non-Markovian
dynamics. Memory kernels or fractional dynamics \cite{carson2025stochastic}
may provide the mechanism.

Further open problems include: characterizing the continuous limit as
$N \to \infty$ and its relation to SDE models
\cite{carson2025statistical}; designing optimal transition kernels for a
given target distribution; and quantifying the acceleration from
non-reversible Markov chain strategies applied to CoT reasoning.
